\section{Frequency and time-domain analysis}

\subsection{Ctrl+F keywords}

\subsubsection{Step response, overshoot, damping, steady-state value, stationaer vaerdi}
\[
\text{DC gain, final value theorem, sinusoidal response, damping ratio, settling time}
\]
Use this section for time-response questions and for checking whether a Bode or controller answer gives plausible dynamics.

\subsection{Steady-state and final value}

\subsubsection{DC gain}
\[
K_{\mathrm{ss}}=G(0)=\lim_{s\to0}G(s)
\]
For a stable system with constant input, \(K_{\mathrm{ss}}\) is the output/input ratio.

\subsubsection{Final value theorem}
\[
\lim_{t\to\infty}y(t)=\lim_{s\to0}sY(s)
\]
All poles of \(sY(s)\) must be in the open left half-plane for this theorem to be valid.

\subsection{Frequency response}

\subsubsection{Evaluate transfer function on the imaginary axis}
\[
G(j\omega)=G(s)\big|_{s=j\omega}
\]
Magnitude and phase describe sinusoidal steady-state behavior.

\subsubsection{Sinusoidal steady-state output}
\[
u(t)=A\sin(\omega t)
\quad\Longrightarrow\quad
y_{\mathrm{ss}}(t)=A|G(j\omega)|\sin\bigl(\omega t+\angle G(j\omega)\bigr)
\]
The magnitude scales the amplitude and the phase shifts the sinusoid.

\subsection{First-order response}

\subsubsection{First-order transfer function and step response}
\[
G(s)=\frac{K}{\tau s+1},\qquad
y(t)=K\left(1-e^{-t/\tau}\right)u(t)
\]
At \(t=\tau\), the response has reached \(63.2\%\) of the total change.

\subsection{Second-order response}

\subsubsection{Standard second-order denominator}
\[
s^2+2\zeta\omega_n s+\omega_n^2
\]
Compare a closed-loop denominator with this form to identify \(\zeta\) and \(\omega_n\).

\subsubsection{Second-order poles}
\[
p_{1,2}=-\zeta\omega_n\pm j\omega_n\sqrt{1-\zeta^2}
\]
Complex poles occur for \(0<\zeta<1\) and produce oscillation and overshoot.

\subsubsection{Percent overshoot}
\[
M_p=100e^{-\zeta\pi/\sqrt{1-\zeta^2}}\%
\]
Use this for underdamped second-order unit-step specifications.

\subsubsection{Settling time approximation}
\[
t_s\approx\frac{4}{\zeta\omega_n}
\]
This is the common \(2\%\) settling-time approximation.

\subsection{Stability from poles}

\subsubsection{Asymptotic stability}
\[
\operatorname{Re}(p_k)<0\quad\text{for all poles }p_k
\]
All poles must lie in the open left half-plane.

\subsubsection{Unstable and marginal pole cases}
\[
\exists p_k:\operatorname{Re}(p_k)>0 \Rightarrow \text{unstable},\qquad
\text{repeated } p_k=0 \Rightarrow \text{unstable}
\]
Simple imaginary-axis poles are marginal only in special cases; repeated poles at the origin are not stable.

\subsection{Exam recipe: overshoot limit from closed-loop denominator}

\textbf{Use when:} the problem asks whether a \(K_P\), \(\zeta\), or closed-loop transfer function satisfies an overshoot bound.

\textbf{Steps:}
\begin{enumerate}
\item Put the denominator in \(s^2+2\zeta\omega_ns+\omega_n^2\) form.
\item Solve for \(\omega_n\) from the constant term.
\item Solve for \(\zeta\) from the \(s\)-coefficient.
\item Insert \(\zeta\) into \(M_p=100e^{-\zeta\pi/\sqrt{1-\zeta^2}}\%\).
\item Compare with the required overshoot.
\end{enumerate}

\textbf{Fast checks:} larger \(\zeta\) means lower overshoot; \(\zeta\geq 1\) gives no oscillatory overshoot in the standard model.

\textbf{Common traps:} applying the formula to a higher-order response without checking dominant second-order behavior.

\subsection{Exam recipe: final value and steady-state output}

\textbf{Use when:} the problem asks for final output, DC gain, or steady-state response.

\textbf{Steps:}
\begin{enumerate}
\item Write \(Y(s)=G(s)U(s)\) or the closed-loop transfer from the diagram.
\item Insert the input transform, often \(1/s\) for a unit step.
\item Check the poles of \(sY(s)\) are stable.
\item Compute \(\lim_{s\to0}sY(s)\).
\end{enumerate}

\textbf{Fast checks:} for a stable step response, the answer should match DC gain times step amplitude.
